\label{chap:methodology}
%% REV (please make a concept that we track that the two studys are always named the same in the thesis: (Formative (needfinding) Study, and Evaluation Study. Again with the first thesis appearence in the glossary)
This chapter sets out the methodological frame within which the research questions posed in Chapter~\ref{chap:introduction} are addressed: the research paradigm that structures the work, the end-to-end process connecting its stages, and the scope, limitations, assumptions, and data-handling practices that bound how the findings should be read.

\section{Research Paradigm}
\label{sec:paradigm}
%% REV (please make a check on the munzner nested model, that all the things are valid we state here. Please also give me your general fit opinion: how well does this model fit the purpose and domain)
This thesis follows Munzner's nested model for visualisation design and validation~\cite{munzner2009}. The nested model organises a visualisation design process into four stratified levels: (1)~domain problem characterisation, (2)~data and task abstraction, (3)~visual encoding and interaction design, and (4)~algorithm design. A defining property of the model, central to how this thesis is structured, is that each level's output is the next level's input. Consequently, an error or unwarranted assumption at an upstream level cascades to every level downstream of it; the model therefore calls for validation at each level individually, rather than only at the end of the pipeline~\cite{munzner2009}.

This thesis applies the model as follows. Level~1 (domain problem characterisation) is addressed in two interconnected parts: through the review of Related Work (Chapter~\ref{chap:related-work}) and through the empirical needfinding study (Chapter~\ref{chap:formative-study}). Together, they define the problem space: the literature provides the theoretical foundation, while the study establishes the empirical reality. Level~2 (data/task abstraction) is addressed by the codebook-based translation of the needfinding results into a features-and-requirements specification (Chapter~\ref{chap:from-themes-to-requirements}). Level~3 (visual encoding and interaction design) and level~4 (algorithm design) are addressed by the Design and Implementation chapters, respectively. Notably, Munzner's own level~2 construct is narrower than what Chapter~\ref{chap:from-themes-to-requirements} produces. The model traditionally calls for generic operations and data types (e.g.,\ filter, sort, compare, applied to a table or a graph), whereas this thesis produces a comprehensive 13-feature requirements specification. This specification does not substitute level~2; rather, the generic operations and data types expected by the model are embedded within this richer document. Munzner explicitly anticipates this kind of adaptation~\cite{munzner2009}; because the nested model acts as an overarching macro-structure rather than a rigid toolkit. It naturally accommodates the integration of domain-specific methods, like standard requirements engineering, where its own prescriptions are narrow.

The two empirical studies in this thesis address distinct, non-overlapping spans of this model. The formative study focuses strictly on the pre-design phase. Operating within the problem space defined in level~1, it produces and establishes the specific requirements for the graph-table hybrid concept. Crucially, it does not evaluate the hybrid concept itself. Whether a combined graph-and-table interface is the optimal approach remains an upstream assumption (stated explicitly in Section~\ref{sec:overall-assumptions}), rather than a hypothesis tested empirically. In contrast, the evaluation study (Chapter~\ref{chap:evaluation}) validates the encoding/interaction and algorithm levels, levels~3 and~4, of the resulting artefact. Crucially, because the nested model posits that downstream validation can expose upstream flaws, the findings from this evaluation inevitably reflect back on that foundational assumption, providing empirical feedback on whether the hybrid graph-table concept was the right approach. This fits the model's call to distinguish contributions by level and state upstream assumptions explicitly~\cite{munzner2009}. Figure~\ref{fig:nested-model-mapping} maps this structure against the process stages introduced in Section~\ref{sec:process-overview}.

\begin{figure}[htbp]
\centering
\begin{tikzpicture}[font=\small, >=stealth, scale=0.85, transform shape]
  % Define box dimensions
  \def\bw{2.1}   % box width
  
  % ---- Level band backgrounds ----
  \fill[subjectblue20, rounded corners=4pt] (-0.15,2.3) rectangle (4.75,3.8);
  \fill[subjectgreen20, rounded corners=4pt] (4.85,2.3) rectangle (7.25,3.8);
  \fill[subjectorange20, rounded corners=4pt] (7.35,2.3) rectangle (9.75,3.8);
  \fill[gray!12, rounded corners=4pt] (9.85,2.3) rectangle (12.25,3.8);

  % ---- Level labels ----
  \node[font=\scriptsize\bfseries, text=subjectblue, align=center, anchor=south] at (2.3, 3.85) {Level 1\\Domain Characterisation};
  \node[font=\scriptsize\bfseries, text=subjectgreen, align=center, anchor=south] at (6.05, 3.85) {Level 2\\Abstraction};
  \node[font=\scriptsize\bfseries, text=subjectorange, align=center, anchor=south] at (8.55, 3.85) {Level 3\\Encoding/\\Interaction};
  \node[font=\scriptsize\bfseries, text=gray!70, align=center, anchor=south] at (11.05, 3.85) {Level 4\\Algorithm};

  % ---- Stage boxes macro ----
  \def\stagebox#1#2#3{
    \fill[white, rounded corners=3pt] (#1,2.5) rectangle (#1+\bw,3.6);
    \draw[#3, line width=1.2pt, rounded corners=3pt] (#1,2.5) rectangle (#1+\bw,3.6);
    \node[align=center, text width=1.9cm, font=\scriptsize] at (#1+1.05,3.05) {#2};
  }

  % ---- Draw the 6 Process Stages (Colored by Level) ----
  \stagebox{0.0}{Related\\Work}{subjectblue}
  \stagebox{2.5}{Formative\\Study}{subjectblue}
  \stagebox{5.0}{Themes $\to$\\Requirements}{subjectgreen}
  \stagebox{7.5}{Design}{subjectorange}
  \stagebox{10.0}{Implementa-\\tion\\\& Testing}{gray!60}
  
  % NEW DISTINCT COLOR for Study 2 (Violet)
  \stagebox{12.5}{Evaluation\\Study}{violet}

  % ---- Gray Process Arrows (Linear Flow) ----
  \draw[->, gray!60, thick] (2.1,3.05) -- (2.5,3.05);
  \draw[->, gray!60, thick] (4.6,3.05) -- (5.0,3.05);
  \draw[->, gray!60, thick] (7.1,3.05) -- (7.5,3.05);
  \draw[->, gray!60, thick] (9.6,3.05) -- (10.0,3.05);
  \draw[->, gray!60, thick] (12.1,3.05) -- (12.5,3.05);

  % ---- Validation spans & feedback loops ----

  % 1. Formative Study (Matches Level 1: Blue)
  \draw[->, dashed, subjectblue, thick, rounded corners=4pt] 
    (3.55, 2.48) -- (3.55, 1.6) -- (6.05, 1.6) -- (6.05, 2.48);
  \node[font=\scriptsize, text=subjectblue, align=center, anchor=north] at (4.8, 1.55) 
    {Formative study\\establishes levels 1--2\\(pre-design)};

  % 2. Evaluation Study (Distinct Color: Violet)
  \draw[dashed, violet, thick, rounded corners=4pt] (13.55, 2.48) -- (13.55, 1.6) -- (8.55, 1.6);
  \draw[->, dashed, violet, thick] (11.05, 1.6) -- (11.05, 2.48);
  \draw[->, dashed, violet, thick] (8.55, 1.6) -- (8.55, 2.48);
  \node[font=\scriptsize, text=violet, align=center, anchor=north] at (11.05, 1.55) 
    {Evaluation study\\validates levels 3--4};

  % 3. Feedback Loop (Matches Study 2: Violet)
  \draw[->, densely dotted, violet, thick, rounded corners=4pt] 
    (13.75, 2.48) -- (13.75, -0.2) -- (3.55, -0.2);
  \node[font=\scriptsize\itshape, text=violet, anchor=north] at (9.275, -0.25) 
    {reflects back on upstream assumptions};

\end{tikzpicture}
\caption{The thesis's process stages mapped onto Munzner's four nested-model levels. The formative study produces and establishes levels~1--2 before design work begins, while the evaluation study validates levels~3--4 after implementation, reflecting its findings back onto the upstream assumptions.}
\label{fig:nested-model-mapping}
\end{figure}

%% REV (i shorten the following chapter so its not that repetetive. But im not happy how it is now. Please find a more structured solution for the block 2 and 3 (1 is fine). Maybe first say what is simiar in both studies, (held with students, think allowed, ...) and then structurally whats the difference. It does not have to be longer (shorter is also fine). While this please also check if the block is factual (are the studies what we claim here?))
\section{Process Overview}
\label{sec:process-overview}
The end-to-end research process operationalises the nested model across six stages: (1)~a review of related work (Chapter~\ref{chap:related-work}); (2)~a formative needfinding study using a non-functional prototype (Chapter~\ref{chap:formative-study}); (3)~translation of findings into a requirements specification (Chapter~\ref{chap:from-themes-to-requirements}); (4)~system design (Chapter~\ref{chap:design}); (5)~software implementation and testing (Chapter~\ref{chap:implementation}); and (6)~an evaluative user study (Chapter~\ref{chap:evaluation}). Because each stage's specific methodology is documented in full within its respective chapter, this section outlines only the umbrella logistics of the two empirical anchor points.

The two empirical studies serve distinct methodological purposes and employ different study designs. The formative study ($N=6$) was an elicitation-focused interview series centred on an interactive, non-functional prototype. Conversely, the evaluation study ($N=11$) was an observational laboratory study of the fully implemented system, designed to capture behavioural usage data. Each evaluation session followed a fixed, within-subjects sequence administered via a purpose-built study-guide application: consent and demographics, a short video tutorial, a first planning scenario and questionnaire, a second scenario and questionnaire, an open-exploration phase, a final user-experience questionnaire, and a semi-structured interview. Participants employed a concurrent think-aloud protocol throughout the working phases, with a moderator present strictly for operational questions. 

To ensure rigorous qualitative analysis, all evaluation sessions were screen- and audio-recorded. The analysis rests on two independent data streams produced from these recordings, a verbatim transcript and a frame-sampled log of interface states, which were subsequently triangulated. Chapter~\ref{chap:evaluation} documents the instruments, the labelling scheme, and the coding procedure in full.



\section{Scope}
\label{sec:overall-scope}

This thesis defines four system-wide scope boundaries, denoted OOS1 through OOS4 %% REV (please introduce a real apprenetion for this)
(Out Of Scope). OOS1 and OOS2 are general enough to be justified independently of any specific finding, and are stated here:

\begin{itemize}
    \item \textbf{OOS1 (no automated planning):} The system guides students in planning; it does not plan on their behalf. This reflects both the intended role of the tool (decision support, not decision automation) and the limited context available to the system to plan responsibly on a student's behalf.
    \item \textbf{OOS2 (programme-level granularity):} Study planning is defined at the programme level throughout this thesis. Course-level planning, scheduling within a single course, deadline management, and appointment tracking are out of scope.
\end{itemize}
%% REV (should we mention them within the list like OOS1 and 2, this looks more consistant. We could still say that we state them later again where they apply and link the Chapter)
OOS3 (no live connection to the institutional CMS/LMS) and OOS4 (no community features) are stated later, in Chapter~\ref{chap:from-themes-to-requirements} (Section~\ref{sec:overall-scope-refined}), alongside the specific excluded findings that justify them.

%% REV (listing this as scope is for sure true, but listing just this for scope feels a little bit wired. Can give me your opinion on this. Also check the thesis for other scope items we could or should mention here)
Two degree programmes at the Faculty of Informatics of TU~Wien are in scope: the Bachelor programme in Computer Science and the Master programme in Software Engineering. Every curriculum rule the system checks, every catalogue entry it serves, and every plan evaluated in either study belongs to one of these two programmes, and the compliance engine holds a separate rule set for each (Chapter~\ref{chap:implementation}). 
The two were chosen because they sit at different stages of a student's path, so that the requirements gathered span both a largely prescribed early curriculum and a largely elective later one.

\section{Limitations and Constraints}
\label{sec:overall-scope-detail}

Beyond the scope boundaries defined in Section~\ref{sec:overall-scope}, four specific methodological limitations bound how the empirical findings in this thesis should be applied:

\begin{itemize}
%% REV (the first limitaion here is also a scop decision earlier. Is this okay)
    \item \textbf{Single-institution scope.} Both the formative-study participants and the initial system's curriculum data are drawn exclusively from TU Wien's Informatics faculty, covering one Bachelor's and one Master's programme. All six formative-study participants were Informatics students; students with different disciplinary backgrounds, or at other institutions, may have different planning needs not captured here.
    \item \textbf{Small sample sizes in both studies.} The formative study involved six participants; the evaluation study analyses eleven participants, which is smaller than would support strong statistical generalisation. Both studies are designed for depth of insight within a bounded population, not for population-level generalisability.
    \item \textbf{Limited direct prior research on hybrid graph-table interfaces specifically.} As Chapter~\ref{chap:related-work} shows, %% REV (is it not none? You could say to the best of our knowledge none...)
    few existing systems combine a graph and a table view for study planning in this way; the evidence motivating the hybrid concept is inferential, drawn from adjacent systems and general visualisation/HCI literature, rather than from a substantial, converging body of research that directly tests hybrid interfaces of this kind against alternatives.
    \item \textbf{Single-researcher qualitative analysis.} The formative study's coding and theme development was conducted by a single researcher, without a second coder or formal inter-coder reliability check. %% REV (this is a arguemnt for the formative study, but what is with the second one? I would eather stay general (limitation of just one researcher) or also have a justification for study 2 in place)
    This is consistent with the choice of codebook TA over coding-reliability TA (Chapter~\ref{chap:formative-study}, Section~\ref{sec:analysis-method}), which does not require multiple coders, but it remains a limitation relative to approaches that do.
\end{itemize}

%% REV (so my stack would be to completly remove the assumptin part here. The one really mentionable assumption is already stated earlier (hyorid graph-table), the others are not important in my opionon. You can check if there are other imporant assumptions we are having, if not delelte it. (or is it  important in research to always have a assumption section?)
\section{Assumptions}
\label{sec:overall-assumptions}

Assumptions, by contrast, are premises this thesis takes as given from the start, before any study was run; they are not things the studies could have avoided with a larger sample or more researchers, since they are not consequences of how the studies were conducted at all, they shape how every later chapter's claims should be read regardless of study design. Two assumptions of this kind are load-bearing enough to state explicitly. The first (as already discussed in Section~\ref{sec:paradigm}) is that the need for a combined graph-and-table interface is motivated by literature (Chapter~\ref{chap:related-work}), not empirically tested within this thesis against a non-hybrid alternative; the formative study validates the requirements derived \emph{within} that concept, not the concept itself. The second concerns what a curriculum is taken to contain. The graph-based literature reviewed in Chapter~\ref{chap:related-work} treats prerequisite sequencing as the planning-relevant structure of a curriculum, and this thesis carries that premise into its design rather than deriving it from the two curricula in scope. Those curricula publish very little sequencing, which is why the graph the evaluation study examines renders the containment hierarchy together with a two-edge prerequisite overlay, and why an advisory about recommended ordering arose only once across twenty-two planning runs (Chapter~\ref{chap:evaluation}; Chapter~\ref{chap:discussion}, Section~\ref{sec:disc-related}). The premise is stated here so that those chapters read as testing it rather than as reporting an omission. The generalisability of TU Wien-specific findings to other institutions, which might otherwise read as a second assumption, is deliberately not repeated here, it is the single-institution-scope limitation above, and stating it under both headings would suggest two distinct concerns where there is only one.


\section{Expected Results}
\label{sec:expected-results}

This thesis produces three core outputs: a ranked features-and-requirements specification grounded in interview evidence (Chapter~\ref{chap:from-themes-to-requirements}); a designed and implemented graph-table hybrid realising it, with every design decision traceable to a requirement, a prototype assumption, common design knowledge, or literature (Chapters~\ref{chap:design}--\ref{chap:implementation}); and empirical findings on how that tool performs against the research questions (Chapter~\ref{chap:evaluation}). 
%% REV (in general good, one question here, are the questionare results a intermidiate state which ended up in the positive findings and usability problems, or are they a seperate output?)
%% REV (the process account is from the recoding, and the usability findings from? Recoring and questionaire?)
cThe third output takes two forms: a process account derived from the interface record and a catalogue of usability problems and positive findings, each carried by a code with a participant-level frequency and mapped to an ISO~9241-11 dimension~\cite{iso9241_11}. Crucially, because the evaluation relies on a fixed scenario order and a small sample size ($N=11$), these two outputs are descriptive rather than causal. They illustrate how students approach planning, but cannot statistically prove that a specific interface component caused a specific behaviour. Given these constraints, Chapter~\ref{chap:discussion} (Section~\ref{sec:disc-generalisability}) examines which of these findings apply exclusively to the graph-table hybrid concept, and which might generalise to other systems.

\section{Data Handling}
\label{sec:ethics-overview}

Both studies followed strict GDPR-compliant consent, recording, and anonymisation protocols. Participants provided explicit consent for data collection and processing (Art.~6(1)(a) GDPR) prior to their sessions, acknowledging their right to withdraw. Sessions were recorded, transcribed, and fully anonymised; no master key linking identities to participant codes (I1--I6; P1--P11) was retained. All raw, identifiable data, including signed forms and internal-only webcam recordings used for frame-sampling, is stored on secure, access-restricted infrastructure and will be permanently deleted upon the conclusion of the thesis process.

To administer the evaluation study, a lightweight, purpose-built web application was developed separately from the main study-planning artefact (Chapter~\ref{chap:evaluation}). This tool presented the scenario briefs, administered post-task questionnaires, and logged responses. Developing a custom application was beneficial for data integrity: it ensured standardised session pacing, enabled precise, unified logging of participant inputs, and eliminated reliance on third-party survey platforms, thereby further securing participant data.

\section{Quality Criteria and Validity Strategy}
\label{sec:validity-overview}

Methodological rigour is systematically addressed across all research stages. The formative study explicitly positions its coding approach within Braun and Clarke's codebook thematic analysis, justifying this choice over alternative paradigms (Section~\ref{sec:analysis-method}). Design and implementation validity is secured through strict traceability: every design decision must trace back to a requirement, prior evidence, or an explicitly stated assumption (Chapter~\ref{chap:design}).

%% REV (the first block here is really nice, for the second i have to say we see that we dont really have a lot of source grounding for our approach in evalstudy. For the first we have a detailed Braun and Clarke justification, for the second, we just do sth. Can we also apply the same reasoning from study 1 to the second, if not we have to come up with references which justify why we analysed the second study like we did)
The evaluation study secures reliability through methodological triangulation, explicitly separating observed behaviour (frame-sampled logs) from self-reported perception (transcripts). Applying a systematic codebook across all participants prevents anecdotal reporting and ensures frequencies reflect genuine prevalence. Where self-reports and observed data conflict, such as participants reporting success despite failing a task brief, the objective observation carries the claim. %% REV (these two limitations are already listed in the limitations section, i think you can cut them out here)
Nonetheless, two analytical limitations remain: severity ratings rely on single-analyst judgement (making problem rankings indicative rather than exact), and quantitative metrics (a domain-adapted acceptance survey and the UEQ) function strictly as descriptive indicators, as the small sample size precludes statistical inference.

%% REV (the following block reads like borderblade, i cannot see content which is actually important, in my opinion this can be deleted, whats your)
Consequently, three principal threats to validity bound the findings in this thesis:
\begin{itemize}
    \item \textbf{Internal validity} is constrained by the evaluation's fixed scenario order, which confounds graph availability with practice effects and scenario topics. Findings are partitioned accordingly to isolate structural effects from sequence effects.
    \item \textbf{Construct validity} is limited by observational proxies: self-reported success measures \emph{perceived} rather than \emph{actual} completion, and interface logs capture panel configuration rather than visual attention.
    \item \textbf{External validity} is limited by the single-institution context, the focus on two specific study programmes, and the small sample sizes in both studies.
\end{itemize}
To ensure transparency, these limitations are not merely aggregated here, but are explicitly reiterated alongside the specific empirical claims they affect.
